\section{Metodología}

La metodología se organizará en cuatro bloques iterativos:

\begin{enumerate}
  \item \textbf{Modelado del sistema.} Definición de estados, parámetros nominales de la carga, restricciones básicas de movimiento y estructura del grafo de comunicación.
  \item \textbf{Diseño de control.} Implementación de una línea base nominal basada en consenso y de una variante adaptativa o robusta para incertidumbre inercial, siguiendo principios de control distribuido y referencias afines en coordinación multiagente \citep{bullo2009distributed,quijano2017population}.
  \item \textbf{Campaña experimental.} Construcción de escenarios reproducibles con configuraciones de carga, perturbaciones y conectividad diferenciadas.
  \item \textbf{Evaluación cuantitativa.} Cálculo de métricas, análisis comparativo y elaboración de figuras y tablas para sustentar conclusiones.
\end{enumerate}

Desde el punto de vista de ingeniería, se adoptará una arquitectura limpia con separación entre modelos de dominio, controladores, simulación, métricas, experimentos, visualización y documentación. Esta organización facilitará trazabilidad, repetibilidad y extensión futura del código.
